Se logró graficar la curva de Franck-Hertz para el mercurio, lo que permitió analizar su comportamiento y observar que la corriente no sigue un patrón lineal. Inicialmente, aumenta con el potencial hasta alcanzar un valor máximo, luego decrece y, finalmente, vuelve a incrementarse. Se determinó que la separación promedio entre los picos es \( \Delta U = \qty{5.2}{\volt} \), con un error relativo de \qty{7.54}{\percent} respecto al valor registrado en la literatura.
